\chapter{Podstawy teoretyczne}
\label{ch:podstawy-teoretyczne}

Rozdział przedstawia podstawy teoretyczne potrzebne do uzasadnienia projektu neuronowego kodeka mowy o bardzo niskiej przepływności bitowej. Zagadnienie to znajduje się na przecięciu klasycznego kodowania mowy, przetwarzania sygnałów, uczenia reprezentacji dyskretnych oraz optymalizacji modeli pod urządzenia o ograniczonych zasobach. Z tego względu w pierwszej kolejności omówiono literaturę dotyczącą kodeków parametrycznych, hybrydowych i neuronowych, a następnie scharakteryzowano aktualny stan technologii w zakresie przepływności poniżej 1 kbps \cite{Rowe2019Codec2,Valin2019LPCNet,Zeghidour2021SoundStream,Defossez2022Encodec}. Dalsze części rozdziału stanowią podstawę do opisu zaproponowanego rozwiązania, ponieważ wyjaśniają właściwości sygnału mowy, cyfrową reprezentację audio, znaczenie przepływności bitowej, podstawy neuronowej kompresji sygnału oraz metryki stosowane do oceny jakości rekonstrukcji.

Warto zaznaczyć, że literatura dotycząca kodeków mowy o ultra-niskiej przepływności nie tworzy jednorodnego pola badawczego. Część prac koncentruje się na klasycznym modelowaniu źródła mowy, część na neuronowym generowaniu sygnału z niewielkiej liczby parametrów, a część na pełnych kodekach uczonych end-to-end. Porównanie tych podejść wymaga ostrożności, ponieważ poszczególne prace stosują różne zbiory danych, różne procedury testów odsłuchowych oraz różne metryki obiektywne. W pracy przyjęto więc perspektywę projektową: istotne są nie tylko raportowane wartości MOS, PESQ lub ViSQOL, ale także możliwość uzyskania niskiej przepływności, stabilność kwantyzacji oraz koszt inferencji na platformie brzegowej.

\section{Analiza literatury zagadnienia}
\label{sec:analiza-literatury}

\subsection{Zakres przeglądu i kryteria porównania}
\label{subsec:zakres-przegladu-literatury}

Przegląd literatury w niniejszej pracy obejmuje publikacje dotyczące ultra-niskobitowej kompresji mowy, neuronowych kodeków audio oraz wdrażania modeli przetwarzania sygnałów na urządzeniach brzegowych. Uwzględniono przede wszystkim prace z lat 2015--2026, ze szczególnym naciskiem na rozwiązania działające przy przepływnościach rzędu pojedynczych kilobitów na sekundę lub niższych. Analizie poddano architekturę kodeka, technikę kwantyzacji, sposób parametryzacji sygnału mowy, metryki jakości oraz wymagania zasobowe. Takie kryteria są istotne, ponieważ w projektowanym systemie sama deklaracja niskiego bitrate'u nie wystarcza; konieczne jest także określenie, czy rozwiązanie może działać w warunkach ograniczonej pamięci, ograniczonej mocy obliczeniowej i małego dopuszczalnego opóźnienia.

Na poziomie technologicznym przegląd wskazuje trzy główne klasy rozwiązań. Pierwszą z nich są kodeki parametryczne, które opisują mowę za pomocą niewielkiego zestawu parametrów, takich jak współczynniki predykcji liniowej, częstotliwość podstawowa i energia. Drugą grupę stanowią kodeki hybrydowe, łączące klasyczne modelowanie mowy z lekkim komponentem neuronowym. Trzecią grupą są kodeki neuronowe end-to-end, w których koder, kwantyzer i dekoder są uczone jako jeden system \cite{Rowe2019Codec2,Valin2019LPCNet,Zeghidour2021SoundStream,Defossez2022Encodec}. Każda z tych rodzin rozwiązań ma inne ograniczenia: kodeki parametryczne są zwykle bardzo lekkie, ale mniej naturalne percepcyjnie; kodeki hybrydowe oferują kompromis między jakością i kosztem; natomiast rozwiązania end-to-end mogą uzyskiwać wysoką jakość, lecz często wymagają większych modeli i kosztowniejszego treningu.

\begin{center}
\centering
\footnotesize
\setlength{\tabcolsep}{3pt}
\begin{tabular}{c c c c c}
\fbox{\parbox{0.19\textwidth}{\centering kodeki\\parametryczne}} &
\(\rightarrow\) &
\fbox{\parbox{0.19\textwidth}{\centering kodeki\\hybrydowe}} &
\(\rightarrow\) &
\fbox{\parbox{0.19\textwidth}{\centering kodeki\\neuronowe end-to-end}} \\
niski koszt & & kompromis jakości i kosztu & & większa elastyczność modelu \\
silny model mowy & & model mowy + sieć & & koder + kwantyzer + dekoder \\
\end{tabular}
\captionof{figure}{Główne rodziny rozwiązań analizowanych w kontekście niskobitowego kodowania mowy.}
\label{fig:rodziny-kodekow}
\end{center}

\subsection{Kodeki parametryczne i hybrydowe}
\label{subsec:kodeki-parametryczne-hybrydowe}

Najniższe przepływności osiągają rozwiązania parametryczne, w których nie rekonstruuje się bezpośrednio pełnego przebiegu audio, lecz przesyła parametry modelu mowy. Przykładem takiego podejścia jest Codec 2, rozwijany jako otwarty kodek niskobitowy i raportowany w wariantach rzędu 700 bit/s oraz 450 bit/s \cite{Rowe2019Codec2,Rowe2021Codec2450}. Tak niska przepływność jest możliwa dzięki silnemu wykorzystaniu założeń o strukturze mowy, jednak wiąże się z ograniczoną naturalnością rekonstrukcji. Zbliżoną rodzinę rozwiązań reprezentuje standard MELPe, opisany w STANAG 4591, który działa w trybach 600, 1200 i 2400 bit/s \cite{NATO2015STANAG4591}. Warto zauważyć, że te rozwiązania pozostają ważnym punktem odniesienia dla pracy, ponieważ pokazują dolną granicę możliwej przepływności, ale nie odpowiadają bezpośrednio na pytanie, jak wykorzystać sieci neuronowe do poprawy jakości odtwarzanego sygnału.

Innym nurtem w tej grupie są metody oparte na predykcji liniowej i efektywnym kodowaniu parametrów. Backstrom i Magi zaproponowali rzadką predykcję liniową jako sposób poprawy efektywności kodowania przy niskich przepływnościach \cite{Backstrom2017SparseLinearPrediction}. W takim podejściu nadal wykorzystuje się klasyczny opis sygnału mowy, ale modyfikuje się sposób reprezentacji i kodowania parametrów. Rozwiązania tego typu są istotne dla projektowanego systemu, ponieważ wskazują, że ograniczenie liczby przesyłanych informacji może być osiągane nie tylko przez zmniejszenie modelu, lecz także przez odpowiedni dobór reprezentacji wejściowej i słownika kodowego.

Kodeki hybrydowe próbują połączyć zalety klasycznego modelowania mowy z większą elastycznością modeli neuronowych. Przykładem jest LPCNet, w którym predykcja liniowa jest wykorzystywana do ograniczenia zadania stojącego przed siecią neuronową, a sam model odpowiada za generowanie lub poprawę sygnału mowy \cite{Valin2019LPCNet}. Valin, Büthe i Vos pokazali również, że rozwiązania oparte na lekkich strukturach rekurencyjnych mogą być optymalizowane pod procesory ARM i pracę w czasie rzeczywistym \cite{Valin2022UltraLowLatencySpeechEnhancement}. Uzupełnieniem tego kierunku jest praca Zhu i Valina, w której opisano optymalizacje LPCNet dla urządzeń embedded \cite{Zhu2021OptimizingLPCNet}. Znaczenie tych prac polega na tym, że nie traktują modelu neuronowego jako pełnego zamiennika klasycznego kodeka, lecz jako komponent poprawiający jakość przy zachowaniu niskiej złożoności.

\subsection{Neuronowe kodeki end-to-end i reprezentacje dyskretne}
\label{subsec:neuronowe-kodeki-reprezentacje-dyskretne}

Rozwój neuronowych kodeków audio znacząco zmienił sposób projektowania systemów kompresji. SoundStream zaproponował architekturę end-to-end obejmującą koder, rezydualną kwantyzację wektorową oraz dekoder generujący sygnał audio \cite{Zeghidour2021SoundStream}. Podobny kierunek rozwinięto w EnCodec, gdzie zastosowano wysokiej jakości neuronową kompresję audio i wielostopniową kwantyzację latentów \cite{Defossez2022Encodec}. Prace te są szczególnie istotne dla niniejszej pracy, ponieważ pokazują, że reprezentacja latentna i RVQ mogą tworzyć efektywny mechanizm zamiany sygnału audio na strumień dyskretnych indeksów. Jednocześnie należy zauważyć, że raportowane warianty tych kodeków zwykle działają przy przepływnościach wyższych niż skrajnie niskie zakresy poniżej 1 kbps, a ich koszt obliczeniowy może być problematyczny dla prostych urządzeń brzegowych.

W literaturze występują również próby poprawy kwantyzacji i jakości rekonstrukcji przez modyfikację struktury słowników. HiFi-Codec wprowadza grupową rezydualną kwantyzację wektorową, która ma zwiększać wierność rekonstrukcji przy kontrolowanej liczbie kodów \cite{Yang2023HiFiCodec}. Z kolei prace oparte na dyskretnych reprezentacjach samonadzorowanych pokazują, że sygnał mowy może być rekonstruowany z tokenów lub cech oderwanych częściowo od pełnej postaci fali \cite{Polyak2021SpeechResynthesis}. SpeechTokenizer rozwija ten kierunek, łącząc reprezentacje semantyczne i akustyczne w kontekście modeli mowy \cite{Zhang2023SpeechTokenizer}. Z perspektywy kodeka o ultra-niskiej przepływności podejścia semantyczne są obiecujące, ponieważ mogą redukować liczbę przesyłanych informacji do treści fonetycznej lub lingwistycznej. Ich ograniczeniem pozostaje jednak koszt ekstrakcji cech i ryzyko utraty informacji potrzebnych do naturalnej rekonstrukcji barwy, prozodii i szczegółów artykulacyjnych.

Osobną grupę stanowią prace, które nie są bezpośrednio pełnymi kodekami mowy, lecz wywarły wpływ na projektowanie neuronowych generatorów audio. WaveNet pokazał, że autoregresyjne modele fali mogą generować wysokiej jakości sygnał audio, choć kosztem znacznej złożoności obliczeniowej \cite{VanDenOord2016WaveNet}. Kleijn i inni zastosowali podejście WaveNet do niskobitowego kodowania mowy \cite{Kleijn2018WaveNetLowRate}. Kankanahalli zaproponował natomiast kodowanie mowy optymalizowane end-to-end z wykorzystaniem głębokich sieci neuronowych \cite{Kankanahalli2018EndToEndSpeechCoding}. W kontekście urządzeń brzegowych istotna jest również praca Morrisona, Xu i Bryana, dotycząca efektywnych neuronowych kodeków audio dla strumieniowania o małym opóźnieniu \cite{Morrison2022EfficientNeuralAudioCodecs}. Uzupełniającym punktem odniesienia są metody sinusoidalne, w których sygnał reprezentowany jest przez składowe harmoniczne i szumowe \cite{Garner2019SinusoidalCoding}. Łącznie prace te pokazują, że wysoka jakość neuronowej rekonstrukcji jest możliwa, ale wymaga szczególnej kontroli złożoności, jeżeli system ma działać na brzegu sieci.

\section{Aktualny stan technologii kodowania mowy i dźwięku poniżej 1 kbps}
\label{sec:stan-technologii-ponizej-1kbps}

\subsection{Rozwiązania rzeczywiście pracujące w zakresie subkilobitowym}
\label{subsec:rozwiazania-subkilobitowe}

W zakresie poniżej 1 kbps najbardziej ugruntowanym punktem odniesienia pozostają rozwiązania parametryczne. Codec 2 w wariancie 700 bit/s oraz eksperymentalny wariant 450 bit/s pokazują, że mowa może być przesyłana przy bardzo małej liczbie bitów, jeśli koder wykorzystuje silne założenia o strukturze sygnału \cite{Rowe2019Codec2,Rowe2021Codec2450}. Podobny charakter ma najniższy tryb MELPe pracujący przy 600 bit/s \cite{NATO2015STANAG4591}. Rozwiązania te pozostają praktyczne pod względem zasobów, jednak ich jakość percepcyjna jest ograniczona. W praktyce oznacza to, że bardzo niski bitrate uzyskuje się kosztem naturalności, szerokości pasma, szczegółów artykulacyjnych i odporności na nietypowe głosy lub warunki akustyczne.

Warto zaznaczyć, że subkilobitowe kodowanie mowy rzadko polega na zachowaniu pełnej informacji akustycznej. Częściej jest to kodowanie informacji wystarczającej do zrozumienia wypowiedzi, a nie do wiernej rekonstrukcji naturalnego brzmienia. Z tego powodu klasyczne kodeki parametryczne są nadal konkurencyjne w najniższych zakresach bitrate'u, mimo że nie odpowiadają współczesnym oczekiwaniom dotyczącym jakości dźwięku. W kontekście niniejszej pracy oznacza to, że neuronowy kodek poniżej 1 kbps musi rozwiązać trudniejszy problem niż kodeki pracujące przy kilku kilobitach na sekundę: powinien utrzymać zrozumiałość mowy, a jednocześnie ograniczyć artefakty typowe dla nadmiernie uproszczonego modelu źródła.

\subsection{Granica między kodekami subkilobitowymi a neuronowymi}
\label{subsec:granica-subkilobit-neural}

Neuronowe kodeki end-to-end, takie jak SoundStream i EnCodec, wyznaczają ważny kierunek technologiczny, ale ich typowe konfiguracje nie były projektowane przede wszystkim dla zakresu poniżej 1 kbps \cite{Zeghidour2021SoundStream,Defossez2022Encodec}. EnCodec raportuje m.in. tryb 1{,}5 kbps, co jest już bliskie zakresowi ultra-niskiemu, lecz nadal przekracza próg 1 kbps i wymaga większych zasobów niż klasyczne kodeki parametryczne \cite{Defossez2022Encodec}. SoundStream i HiFi-Codec pokazują natomiast, że RVQ oraz warianty grupowej kwantyzacji są skuteczne przy budowie dyskretnych reprezentacji audio, ale przeniesienie tych metod do zakresu subkilobitowego wymaga znacznego zmniejszenia liczby przesyłanych indeksów lub rozmiarów słowników \cite{Zeghidour2021SoundStream,Yang2023HiFiCodec}. Nowsze prace neuronowe zaczynają schodzić do przepływności subkilobitowych; przykładowo transformerowy kodek z wąskim gardłem FSQ raportuje warianty 400 i 700 bit/s \cite{Parker2025ScalingTransformers}. Wyniki te są istotne dla kierunku badań, lecz należy je traktować ostrożnie w kontekście wdrożeniowym, ponieważ rozwiązania tego typu często wymagają dużych modeli i nie są jeszcze tak ugruntowane praktycznie jak klasyczne kodeki mowy.

Istotną luką technologiczną pozostaje więc zakres między klasycznymi kodekami parametrycznymi a dużymi neuronowymi kodekami audio. Kodeki parametryczne działają poniżej 1 kbps, ale oferują ograniczoną jakość. Kodeki neuronowe często oferują wyższą jakość w raportowanych warunkach, ale zwykle przy większym bitrate lub wyższym koszcie obliczeniowym. Podejścia hybrydowe, takie jak LPCNet, stanowią próbę wypełnienia tej luki, ponieważ wykorzystują parametryczny opis mowy do ograniczenia zadania stawianego sieci neuronowej \cite{Valin2019LPCNet,Zhu2021OptimizingLPCNet}. Wciąż jednak większość praktycznych wariantów hybrydowych raportowana jest dla przepływności około 1{,}6 kbps lub wyższej, co oznacza, że zejście poniżej 1 kbps wymaga dalszej redukcji informacji przesyłanej przez kanał.

\subsection{Podejścia semantyczne i ograniczenia urządzeń brzegowych}
\label{subsec:podejscia-semantyczne-edge}

Jednym z najbardziej interesujących kierunków dla przyszłych kodeków subkilobitowych są reprezentacje semantyczne. SpeechTokenizer oraz prace dotyczące rekonstrukcji mowy z dyskretnych reprezentacji samonadzorowanych sugerują, że część informacji akustycznej można zastąpić tokenami opisującymi treść lub strukturę fonetyczną wypowiedzi \cite{Zhang2023SpeechTokenizer,Polyak2021SpeechResynthesis}. Takie podejście może być korzystne przy bardzo niskiej przepływności, ponieważ liczba tokenów semantycznych może być mniejsza niż liczba informacji potrzebnych do pełnej rekonstrukcji fali. Jednocześnie metoda ta wprowadza nowe ograniczenia: model ekstrakcji reprezentacji semantycznej bywa duży, a rekonstrukcja naturalnego brzmienia wymaga dodatkowych informacji akustycznych, których nie da się łatwo zmieścić przy bardzo niskim limicie przepływności.

W przypadku przetwarzania brzegowego kluczowe znaczenie mają opóźnienie, pamięć i możliwość pracy bez akceleratora GPU. Prace dotyczące optymalizacji LPCNet i modeli mowy na procesorach ARM pokazują, że lekkie architektury oraz wykorzystanie instrukcji SIMD/NEON mogą istotnie zmniejszyć koszt inferencji \cite{Valin2022UltraLowLatencySpeechEnhancement,Zhu2021OptimizingLPCNet}. Z kolei badania nad efektywnymi neuronowymi kodekami audio dla strumieniowania o małym opóźnieniu wskazują, że architektura modelu musi być projektowana razem z wymaganiami wdrożeniowymi, a nie dopiero optymalizowana po zakończeniu treningu \cite{Morrison2022EfficientNeuralAudioCodecs}. Dla projektowanego kodeka oznacza to konieczność kontrolowania nie tylko jakości rekonstrukcji, ale także długości okna kodowania, liczby operacji w koderze i dekoderze oraz rozmiaru pamięci potrzebnej do przechowywania parametrów modelu.

Podsumowując, aktualny stan technologii pokazuje, że poniżej 1 kbps najbardziej ugruntowanym punktem odniesienia pozostają kodeki parametryczne, natomiast rozwiązania neuronowe nadal wymagają kompromisów między jakością, złożonością i dojrzałością wdrożeniową. Najbardziej uzasadnionym kierunkiem dla pracy jest więc projekt inspirowany neuronowymi kodekami end-to-end, ale silnie ograniczony pod względem przepływności i złożoności. W praktyce oznacza to wykorzystanie kodera generującego bardzo zwartą reprezentację latentną, kwantyzacji RVQ lub jej uproszczonego wariantu oraz dekodera zaprojektowanego tak, aby przy małej liczbie bitów rekonstruował przede wszystkim informację istotną percepcyjnie dla sygnału mowy.

\begin{center}
\centering
\footnotesize
\setlength{\tabcolsep}{2pt}
\begin{tabular}{c c c c c}
\fbox{\parbox{0.20\textwidth}{\centering kodeki\\subkilobitowe\\parametryczne}} &
\(\longleftrightarrow\) &
\fbox{\parbox{0.20\textwidth}{\centering luka\\projektowa}} &
\(\longleftrightarrow\) &
\fbox{\parbox{0.20\textwidth}{\centering kodeki\\neuronowe\\wyższej jakości}} \\
poniżej 1 kbps & & cel niniejszej pracy & & zwykle powyżej 1 kbps \\
niski koszt, niższa naturalność & & jakość przy małym koszcie & & lepsza jakość, większy model \\
\end{tabular}
\captionof{figure}{Uproszczona interpretacja luki technologicznej między kodekami parametrycznymi a neuronowymi.}
\label{fig:luka-technologiczna}
\end{center}

\section{Charakterystyka sygnału mowy i dźwięku}
\label{sec:charakterystyka-sygnalu-mowy}

Sygnał mowy jest szczególnym przypadkiem sygnału dźwiękowego, ponieważ łączy właściwości fizyczne fali akustycznej z informacją językową i percepcyjną. Z punktu widzenia kompresji nie jest konieczne zachowanie każdego szczegółu przebiegu próbkowego, lecz tych cech, które pozwalają odbiorcy rozpoznać treść wypowiedzi, intonację oraz podstawowe właściwości głosu. Właśnie dlatego kodeki mowy różnią się od ogólnych kodeków audio: mogą silniej wykorzystywać wiedzę o sposobie wytwarzania mowy, ale jednocześnie są bardziej wrażliwe na utratę informacji fonetycznej, takich jak przejścia między głoskami, struktura formantowa i obecność składowych szumowych \cite{RabinerSchafer2011SpeechDSP,Stevens1998AcousticPhonetics}.

\subsection{Model źródło-filtr i struktura artykulacyjna}
\label{subsec:model-zrodlo-filtr}

Jednym z podstawowych opisów sygnału mowy jest model źródło-filtr, zgodnie z którym sygnał powstaje przez pobudzenie układu artykulacyjnego, a następnie jego przefiltrowanie przez trakt głosowy \cite{Fant1960AcousticTheory,Stevens1998AcousticPhonetics}. W przypadku głosek dźwięcznych źródłem jest zwykle quasi-okresowe pobudzenie związane z pracą fałdów głosowych. W przypadku głosek bezdźwięcznych większe znaczenie ma pobudzenie szumowe, wynikające z turbulentnego przepływu powietrza. Trakt głosowy, obejmujący między innymi jamę ustną i nosową, kształtuje widmo sygnału i prowadzi do powstawania formantów, czyli pasm częstotliwości szczególnie ważnych dla rozróżniania samogłosek oraz części spółgłosek.

Taki opis jest użyteczny w projektowaniu kodeków, ponieważ rozdziela informację o pobudzeniu od informacji o obwiedni widmowej. Klasyczne kodeki parametryczne wykorzystują to rozdzielenie bezpośrednio, przesyłając parametry związane z częstotliwością podstawową, energią i kształtem widma. W kodeku neuronowym podział ten nie musi być jawnie narzucony, ale nadal wpływa na to, czego model powinien nauczyć się w reprezentacji latentnej. Jeżeli koder utraci informację o dźwięczności albo o szybkich zmianach artykulacyjnych, dekoder może wygenerować sygnał brzmiący płynnie, lecz mniej zrozumiały.

\begin{center}
\centering
\footnotesize
\setlength{\tabcolsep}{3pt}
\begin{tabular}{c c c c c}
\fbox{\parbox{0.17\textwidth}{\centering pobudzenie\\okresowe lub szumowe}} &
\(\rightarrow\) &
\fbox{\parbox{0.17\textwidth}{\centering trakt\\głosowy}} &
\(\rightarrow\) &
\fbox{\parbox{0.17\textwidth}{\centering fala\\mowy}} \\
 & & \(\downarrow\) & & \(\downarrow\) \\
 & & \fbox{\parbox{0.17\textwidth}{\centering formanty\\i obwiednia widma}} &
\(\rightarrow\) &
\fbox{\parbox{0.17\textwidth}{\centering cechy\\dla kodeka}} \\
\end{tabular}
\captionof{figure}{Uproszczony model powstawania sygnału mowy i jego znaczenie dla kodowania.}
\label{fig:model-zrodlo-filtr}
\end{center}

\subsection{Cechy czasowo-częstotliwościowe istotne dla kodowania}
\label{subsec:cechy-czasowo-czestotliwosciowe}

Mowa jest sygnałem niestacjonarnym, dlatego jej właściwości zmieniają się w czasie. Krótkie fragmenty mogą być jednak traktowane jako lokalnie stabilne, co uzasadnia analizę ramkową stosowaną w wielu metodach przetwarzania mowy \cite{RabinerSchafer2011SpeechDSP}. W każdej ramce istotne są inne składniki: dla samogłosek ważna jest struktura harmoniczna i formantowa, dla spółgłosek zwartych krótkie przejścia energii, a dla głosek szczelinowych rozkład szumu w wyższych pasmach. Przy bardzo niskiej przepływności problemem nie jest więc wyłącznie redukcja liczby bitów, lecz także decyzja, które zmiany sygnału muszą zostać zachowane z wysoką rozdzielczością czasową.

W praktyce oznacza to, że koder powinien dobrze reprezentować zarówno dłuższe zależności, takie jak intonacja i rytm wypowiedzi, jak i krótkie zdarzenia, takie jak początek spółgłoski lub nagła zmiana energii. Zbyt silne wygładzenie reprezentacji może poprawić stabilność rekonstrukcji, ale jednocześnie pogorszyć artykulację. Z kolei reprezentacja zbyt szczegółowa zwiększa przepływność lub wymaga większego słownika kodowego. Taki kompromis jest szczególnie widoczny w systemach z kwantyzacją wektorową, gdzie każda ramka latentna musi zostać zastąpiona jednym lub kilkoma indeksami.

\begin{center}
\footnotesize
\setlength{\tabcolsep}{4pt}
\begin{tabularx}{0.96\textwidth}{>{\raggedright\arraybackslash}p{0.22\textwidth} >{\raggedright\arraybackslash}X >{\raggedright\arraybackslash}X}
\toprule
\textbf{Cecha sygnału} & \textbf{Znaczenie dla odbioru mowy} & \textbf{Konsekwencja dla kodeka} \\
\midrule
Częstotliwość podstawowa i harmoniczne & Wpływają na wysokość głosu, dźwięczność oraz naturalność wypowiedzi. & Reprezentacja musi zachować informację o okresowości, szczególnie dla głosek dźwięcznych. \\
Formanty i obwiednia widmowa & Ułatwiają rozróżnianie samogłosek i części cech artykulacyjnych. & Zbyt uboga reprezentacja widmowa może prowadzić do utraty zrozumiałości. \\
Fragmenty szumowe & Są istotne dla głosek bezdźwięcznych i szczelinowych. & Dekoder powinien odtwarzać energię w wyższych pasmach bez nadmiernego wygładzania. \\
Przejścia czasowe & Opisują granice fonemów, zwarcia i szybkie zmiany artykulacji. & Zbyt duża redukcja czasowa może usuwać krótkie, ale ważne zdarzenia. \\
\bottomrule
\end{tabularx}
\captionof{table}{Wybrane cechy sygnału mowy istotne z punktu widzenia kompresji.}
\label{tab:cechy-sygnalu-mowy}
\end{center}

\subsection{Dynamika, percepcja i artefakty rekonstrukcji}
\label{subsec:dynamika-percepcja-artefakty}

Ocena jakości mowy nie jest tożsama z pomiarem błędu próbkowego. Dwa sygnały mogą różnić się punktowo w dziedzinie czasu, a mimo to być odbierane podobnie, jeżeli zachowują zbliżoną strukturę widmową i rytmiczną. Z drugiej strony niewielkie zniekształcenie w obszarze ważnym percepcyjnie może być łatwo zauważalne. Psychoakustyka wskazuje, że odbiór dźwięku zależy między innymi od maskowania, głośności, rozdzielczości częstotliwościowej słuchu i wrażliwości na zmiany czasowe \cite{ZwickerFastl1999Psychoacoustics}. W przypadku kodeka mowy oznacza to, że minimalizacja prostego błędu rekonstrukcji nie zawsze prowadzi do najlepszego efektu odsłuchowego.

Najczęstsze artefakty przy silnej kompresji mowy wynikają z utraty szczegółów w czasie i częstotliwości. Mogą one przyjmować postać metalicznego brzmienia, nadmiernego wygładzenia, szumu tła, zniekształcenia spółgłosek lub niestabilnej barwy głosu. W kodekach neuronowych dodatkowym problemem może być generowanie sygnału pozornie naturalnego, ale częściowo niezgodnego z treścią wejściową. Jest to szczególnie istotne w zakresie poniżej 1 kbps, ponieważ ograniczona liczba bitów zwiększa presję na dekoder, który musi odtwarzać szczegóły na podstawie silnie ograniczonej reprezentacji.

Z perspektywy niniejszej pracy najważniejsze jest więc traktowanie sygnału mowy jako struktury wielopoziomowej. Koder powinien przenosić informacje o treści fonetycznej, dźwięczności, energii, obwiedni widmowej i rytmie, a dekoder powinien rekonstruować przebieg w sposób spójny percepcyjnie. Nie oznacza to pełnej wierności akustycznej, ponieważ przy przepływności poniżej 1 kbps jest ona mało realistyczna. Oznacza natomiast konieczność takiego zaprojektowania reprezentacji latentnej i funkcji strat, aby ograniczona liczba kodów była przeznaczana przede wszystkim na informacje istotne dla zrozumiałości oraz naturalności mowy.

\section{Podstawy cyfrowej reprezentacji sygnału audio}
\label{sec:cyfrowa-reprezentacja-audio}

Cyfrowe przetwarzanie dźwięku wymaga zastąpienia ciągłej fali akustycznej skończoną sekwencją liczb. Z punktu widzenia kodeka jest to pierwszy etap redukcji opisu sygnału: dźwięk zostaje ograniczony przez częstotliwość próbkowania, rozdzielczość amplitudy oraz sposób zapisu próbek. Dalsze etapy, takie jak analiza widmowa, ekstrakcja cech melowych lub kodowanie przez sieć neuronową, działają już na tej cyfrowej reprezentacji. Dlatego wybór sposobu reprezentacji nie jest wyłącznie kwestią techniczną, lecz wpływa na informację dostępną dla kodera, funkcje strat oraz ocenę jakości rekonstrukcji.

\subsection{Próbkowanie, kwantyzacja i zapis PCM}
\label{subsec:probkowanie-kwantyzacja-pcm}

Próbkowanie polega na pomiarze wartości sygnału analogowego w równych odstępach czasu. Jeżeli ciągły sygnał audio oznaczono jako \(x_a(t)\), a częstotliwość próbkowania jako \(F_s\), to odpowiadający mu sygnał dyskretny można zapisać w postaci:

\begin{equation}
    x[n] = x_a(nT_s), \qquad T_s = \frac{1}{F_s}.
\end{equation}

Z twierdzenia o próbkowaniu wynika, że sygnał ograniczony pasmowo może zostać odtworzony bez utraty informacji, jeżeli częstotliwość próbkowania jest co najmniej dwukrotnie większa od najwyższej składowej obecnej w sygnale \cite{RabinerSchafer2011SpeechDSP}. W praktyce oznacza to, że sygnał próbkowany z częstotliwością 16 kHz opisuje pasmo do około 8 kHz, co odpowiada mowie szerokopasmowej i obejmuje znacznie więcej informacji niż klasyczna mowa telefoniczna wąskopasmowa. Nie obejmuje jednak najwyższych składowych obecnych w nagraniach muzycznych lub w pełnopasmowym audio. W pracy dotyczącej kodeka mowy taki wybór jest uzasadniony, ponieważ zmniejsza liczbę próbek do przetworzenia i obniża koszt obliczeniowy, a jednocześnie zachowuje znaczną część informacji potrzebnej do rozumienia wypowiedzi.

Drugim etapem cyfryzacji jest kwantyzacja amplitudy. W reprezentacji PCM każda próbka jest zapisywana jako liczba o ustalonej liczbie bitów, najczęściej po przeskalowaniu amplitudy do ograniczonego zakresu. Dla sygnału monofonicznego nominalną przepływność nieskompresowanego zapisu PCM można określić zależnością:

\begin{equation}
    R_{\mathrm{PCM}} = F_s \cdot B,
\end{equation}

gdzie \(B\) oznacza liczbę bitów przypadających na jedną próbkę. Dla sygnału 16 kHz i rozdzielczości 16 bitów daje to 256 kbps. Wartość ta stanowi punkt odniesienia dla dalszej kompresji: kodek pracujący poniżej 1 kbps musi więc zmniejszyć opis sygnału o ponad dwa rzędy wielkości. Tak duża redukcja nie może polegać jedynie na prostym zmniejszeniu liczby bitów próbek, ponieważ prowadziłoby to do silnego szumu kwantyzacji. Konieczne jest utworzenie reprezentacji pośredniej, która przenosi bardziej użyteczne informacje niż pojedyncze wartości amplitudy.

\begin{center}
\footnotesize
\setlength{\tabcolsep}{3pt}
\begin{tabular}{c c c c c}
\fbox{\parbox{0.16\textwidth}{\centering fala\\akustyczna}} &
\(\rightarrow\) &
\fbox{\parbox{0.17\textwidth}{\centering próbki\\PCM}} &
\(\rightarrow\) &
\fbox{\parbox{0.18\textwidth}{\centering ramki\\czasowe}} \\
& & & & \(\downarrow\) \\
\fbox{\parbox{0.17\textwidth}{\centering cechy\\melowe}} &
\(\leftarrow\) &
\fbox{\parbox{0.17\textwidth}{\centering widmo\\STFT}} &
\(\leftarrow\) &
\fbox{\parbox{0.18\textwidth}{\centering okienkowanie\\sygnału}} \\
\end{tabular}
\captionof{figure}{Typowe przejście od fali akustycznej do cyfrowych reprezentacji wykorzystywanych w analizie i uczeniu kodeka.}
\label{fig:cyfrowa-reprezentacja-audio}
\end{center}

\subsection{Reprezentacja fali i analiza STFT}
\label{subsec:reprezentacja-fali-stft}

Najbardziej bezpośrednią reprezentacją cyfrowego audio jest przebieg czasowy, czyli sekwencja próbek. Taki zapis zachowuje informację o fazie, amplitudzie i dokładnym położeniu zdarzeń w czasie, dlatego jest naturalnym wejściem oraz wyjściem dla kodeków rekonstruujących falę dźwiękową. Ma jednak istotne ograniczenie: pojedyncze próbki nie opisują wprost struktury harmonicznej, formantów ani rozkładu energii w pasmach częstotliwości. Dla modelu uczonego end-to-end oznacza to konieczność samodzielnego wyodrębnienia cech, które w klasycznych metodach przetwarzania sygnałów były obliczane jawnie.

Z tego powodu w analizie audio często stosuje się krótkoczasową transformatę Fouriera. STFT dzieli sygnał na zachodzące na siebie ramki, mnoży je przez okno analizy, a następnie oblicza widmo dla każdej ramki. Dla okna \(w[n]\), kroku ramki \(H\) oraz długości transformaty \(N\) można zapisać:

\begin{equation}
    X(m,k) = \sum_{n=0}^{N-1} x[n + mH] w[n] e^{-j2\pi kn/N}.
\end{equation}

Wynikiem jest opis czasowo-częstotliwościowy, który pokazuje, jak energia sygnału rozkłada się w kolejnych pasmach i momentach czasu. W mowie jest to szczególnie przydatne, ponieważ sygnał nie jest stacjonarny, ale krótkie fragmenty mogą być analizowane jako lokalnie stabilne \cite{RabinerSchafer2011SpeechDSP}. Długość okna decyduje o kompromisie między rozdzielczością czasową i częstotliwościową: krótkie okno lepiej lokalizuje szybkie przejścia, natomiast dłuższe okno dokładniej rozdziela składowe harmoniczne. W kodekach neuronowych STFT bywa wykorzystywana nie tylko do wizualizacji, lecz także w funkcjach strat, ponieważ błąd w dziedzinie widmowej często lepiej odpowiada artefaktom słyszalnym niż sam błąd próbkowy.

Należy jednak zauważyć, że widmo STFT nie jest neutralnym zamiennikiem przebiegu czasowego. Moduł transformaty dobrze opisuje rozkład energii, lecz faza pozostaje trudniejsza do modelowania, a jej utrata lub niedokładna rekonstrukcja może prowadzić do zniekształceń. Dlatego systemy rekonstruujące audio zwykle traktują przebieg czasowy jako docelowy sygnał wyjściowy, a reprezentacje widmowe jako narzędzie pomocnicze w ocenie i uczeniu. Takie podejście jest spójne z celem kodeka mowy: model powinien odtwarzać falę dźwiękową, ale jego trening może korzystać z cech, które lepiej opisują strukturę percepcyjną sygnału.

\subsection{Skala melowa i konsekwencje pracy z sygnałem 16 kHz}
\label{subsec:skala-melowa-16khz}

Skala melowa jest nieliniową skalą częstotliwości, która przybliża sposób, w jaki człowiek postrzega różnice wysokości dźwięku. Dla niskich częstotliwości rozdzielczość percepcyjna jest większa, natomiast dla wyższych pasm kolejne różnice częstotliwości są odbierane mniej szczegółowo. Jednym z często stosowanych przybliżeń jest zależność:

\begin{equation}
    m(f) = 2595 \log_{10}\left(1 + \frac{f}{700}\right),
\end{equation}

gdzie \(f\) oznacza częstotliwość w hercach. W praktyce spektrogram melowy powstaje przez przemnożenie widma STFT przez bank filtrów rozmieszczonych według tej skali. Takie przekształcenie zmniejsza wymiar reprezentacji i podkreśla pasma ważne percepcyjnie, dlatego jest powszechnie stosowane w rozpoznawaniu mowy, syntezie mowy oraz funkcjach strat dla modeli audio \cite{RabinerSchafer2011SpeechDSP,ZwickerFastl1999Psychoacoustics}.

W kontekście kompresji ultra-niskobitowej reprezentacje melowe mają dwojakie znaczenie. Z jednej strony pozwalają opisywać sygnał w sposób bardziej zgodny z percepcją niż równomierna skala liniowa. Może to ułatwiać ocenę, czy rekonstrukcja zachowuje barwę i energię w istotnych pasmach. Z drugiej strony samo przejście do skali melowej usuwa część szczegółów, szczególnie w wysokich częstotliwościach, dlatego nie powinno być traktowane jako pełny opis sygnału. W kodeku, którego zadaniem jest odtworzenie fali, cechy melowe są więc bardziej naturalne jako składnik funkcji straty lub narzędzie diagnostyczne niż jako jedyna informacja przesyłana przez kanał.

Dobór częstotliwości 16 kHz wpływa na wszystkie opisane reprezentacje. W dziedzinie czasu oznacza 16 000 próbek na sekundę, co bezpośrednio określa koszt przetwarzania przez koder i dekoder. W analizie STFT ogranicza najwyższe dostępne pasmo do 8 kHz, a w spektrogramie melowym determinuje zakres banku filtrów. Dla mowy jest to kompromis korzystny z punktu widzenia urządzeń brzegowych: redukuje liczbę operacji względem sygnału 44{,}1 kHz lub 48 kHz, a jednocześnie zachowuje pasmo wystarczające do oceny zrozumiałości i większości cech artykulacyjnych. Ograniczeniem pozostaje utrata składowych powyżej 8 kHz, które mogą wpływać na naturalność, ostrość spółgłosek szczelinowych i wrażenie przestrzenności. W niniejszej pracy przyjęcie 16 kHz należy więc traktować jako świadomy kompromis między jakością, kosztem obliczeniowym i docelowym zastosowaniem na urządzeniach brzegowych.

\section{Przepływność bitowa i jej wpływ na jakość dźwięku}
\label{sec:przeplywnosc-jakosc}

Przepływność bitowa określa liczbę bitów potrzebnych do opisania jednej sekundy sygnału. W kompresji audio jest to podstawowa miara kosztu transmisji lub zapisu, ale nie opisuje samodzielnie jakości rekonstrukcji. Dwa kodeki o tej samej przepływności mogą uzyskiwać różną jakość, jeżeli inaczej wykorzystują dostępną liczbę bitów: jeden może przeznaczać więcej informacji na strukturę widmową, inny na dokładniejsze odwzorowanie czasu, a jeszcze inny na parametry modelu mowy. Z tego względu przepływność należy interpretować łącznie z rodzajem reprezentacji, opóźnieniem, złożonością obliczeniową i sposobem oceny jakości.

\subsection{Nominalna i rzeczywista przepływność strumienia}
\label{subsec:nominalna-rzeczywista-przeplywnosc}

W najprostszym przypadku przepływność jest iloczynem liczby jednostek opisu generowanych w ciągu sekundy oraz liczby bitów przypadających na jedną jednostkę. Dla zapisu PCM jednostką jest próbka, dlatego przepływność zależy bezpośrednio od częstotliwości próbkowania i liczby bitów próbki. W kodekach ramkowych jednostką może być natomiast ramka parametrów, wektor cech, indeks słownika lub grupa indeksów opisujących krótki fragment sygnału. Ogólną zależność można zapisać jako:

\begin{equation}
    R = f_{\mathrm{ram}} \cdot b_{\mathrm{ram}},
\end{equation}

gdzie \(f_{\mathrm{ram}}\) oznacza liczbę ramek na sekundę, a \(b_{\mathrm{ram}}\) liczbę bitów przypadających na jedną ramkę. Taki zapis jest użyteczny, ponieważ oddziela decyzje czasowe od decyzji kwantyzacyjnych. Zmniejszenie liczby ramek na sekundę redukuje przepływność, ale pogarsza zdolność kodeka do opisu szybkich zmian mowy. Zmniejszenie liczby bitów w ramce ogranicza pojemność reprezentacji, lecz nie musi zmieniać opóźnienia ani rytmu pracy modelu.

W praktyce należy rozróżnić przepływność nominalną i rzeczywistą. Przepływność nominalna dotyczy samego ładunku informacyjnego, na przykład indeksów wybieranych przez kwantyzer. Nie uwzględnia ona narzutu kontenera, nagłówków, synchronizacji ramek, ewentualnych znaczników długości ani mechanizmów korekcji błędów. Przepływność rzeczywista może być więc wyższa, szczególnie w krótkich pakietach lub systemach strumieniowych. Z drugiej strony zastosowanie kodowania entropijnego może zmniejszyć średnią liczbę bitów, jeżeli rozkład symboli jest nierównomierny. W kodekach przeznaczonych dla urządzeń brzegowych taka optymalizacja wymaga jednak ostrożności, ponieważ może zwiększać złożoność dekodowania, opóźnienie i wrażliwość na błędy transmisji.

\subsection{Przepływność reprezentacji latentnej}
\label{subsec:przeplywnosc-reprezentacji-latentnej}

W neuronowym kodeku z dyskretnym wąskim gardłem przepływność zależy od częstotliwości ramek latentnych oraz liczby bitów potrzebnych do zapisania indeksów kwantyzacji. Jeżeli koder tworzy \(f_{\mathrm{lat}}\) ramek latentnych na sekundę, a każda ramka jest opisana przez kilka indeksów słownikowych, to liczba bitów w ramce wynika z rozmiarów słowników. Dla \(N_q\) słowników o rozmiarach \(K_i\) można przyjąć:

\begin{equation}
    b_{\mathrm{ram}} = \sum_{i=1}^{N_q} \lceil \log_2 K_i \rceil.
\end{equation}

Zapis z funkcją sufitu \(\lceil \cdot \rceil\) odpowiada stałoszerokiemu kodowaniu indeksów, w którym każdy indeks musi zajmować całkowitą liczbę bitów. Wartość \(\log_2 K_i\) opisuje natomiast nominalną informację związaną z wyborem jednego kodu z \(K_i\) elementów i jest naturalna przy analizie teoretycznej lub przy założeniu idealnego kodowania entropijnego. Dla słownika o \(32\) kodach obie wartości są takie same, ponieważ \(\log_2 32 = \lceil \log_2 32 \rceil = 5\).

Zależność ta pokazuje, że bardzo mała zmiana parametrów może znacząco wpłynąć na końcowy bitrate. Zwiększenie liczby słowników poprawia pojemność reprezentacji, ponieważ każda ramka może zostać opisana dokładniej, ale koszt rośnie liniowo względem liczby etapów kwantyzacji. Zwiększenie rozmiaru słownika daje większy wybór kodów, lecz podnosi liczbę bitów potrzebnych na indeks. Z kolei zmniejszenie częstotliwości ramek latentnych pozwala ograniczyć bitrate bez zmiany słowników, ale może utrudniać rekonstrukcję przejść fonetycznych, spółgłosek zwartych i krótkich zmian energii.

Dla zilustrowania tej zależności można rozważyć sygnał próbkowany z częstotliwością \(16\ \mathrm{kHz}\), którego koder redukuje rozdzielczość czasową \(256\)-krotnie. Otrzymuje się wtedy \(62{,}5\) ramek latentnych na sekundę. Jeżeli każda ramka jest opisana przez trzy indeksy RVQ, a każdy słownik ma \(32\) kody, to jeden indeks wymaga nominalnie \(5\) bitów, a cała ramka \(15\) bitów. Przepływność samego strumienia indeksów wynosi wówczas:

\begin{equation}
    R = 62{,}5 \cdot 15 = 937{,}5\ \mathrm{bit/s}.
\end{equation}

Taki przykład pokazuje, dlaczego zakres poniżej 1 kbps jest szczególnie wymagający. Przy tej przepływności jedna sekunda mowy jest opisywana liczbą bitów porównywalną z około setką znaków prostego tekstu. Kodek nie może więc zachowywać pełnej informacji akustycznej, lecz musi nauczyć się przenosić informacje najbardziej użyteczne dla rekonstrukcji mowy. Jednocześnie nominalna liczba bitów nie gwarantuje jeszcze efektywnego wykorzystania kanału. Jeżeli słownik kodowy zapada się i używana jest tylko niewielka część kodów, rzeczywista pojemność reprezentacji maleje, mimo że formalna przepływność pozostaje bez zmian.

\begin{center}
\footnotesize
\setlength{\tabcolsep}{4pt}
\begin{tabularx}{0.96\textwidth}{>{\raggedright\arraybackslash}p{0.25\textwidth} >{\raggedright\arraybackslash}X >{\raggedright\arraybackslash}X}
\toprule
\textbf{Decyzja projektowa} & \textbf{Wpływ na przepływność} & \textbf{Możliwy wpływ na jakość} \\
\midrule
Zmniejszenie liczby ramek latentnych & Redukuje liczbę indeksów przesyłanych w sekundzie. & Może usuwać krótkie zdarzenia i pogarszać artykulację. \\
Zmniejszenie liczby słowników RVQ & Ogranicza liczbę indeksów w każdej ramce. & Może zwiększać błąd kwantyzacji i wygładzać rekonstrukcję. \\
Zmniejszenie rozmiaru słownika & Zmniejsza liczbę bitów potrzebnych na pojedynczy indeks. & Może ograniczać różnorodność kodów i utrudniać opis nietypowych fragmentów mowy. \\
Kodowanie entropijne indeksów & Może obniżyć średnią przepływność przy nierównym rozkładzie symboli. & Może zwiększać złożoność, opóźnienie i zależność od stabilnego rozkładu kodów. \\
\bottomrule
\end{tabularx}
\captionof{table}{Wpływ wybranych decyzji projektowych na przepływność i jakość rekonstrukcji.}
\label{tab:przeplywnosc-decyzje-projektowe}
\end{center}

\subsection{Zależność między bitrate'em a jakością rekonstrukcji}
\label{subsec:bitrate-jakosc-rekonstrukcji}

Wpływ przepływności na jakość dźwięku nie jest liniowy. Przy wysokich przepływnościach dodatkowe bity mogą poprawiać szczegóły, ale różnica percepcyjna bywa niewielka. Przy bardzo niskich przepływnościach każdy bit ma większe znaczenie, ponieważ decyduje o tym, czy model zachowa informację o strukturze fonetycznej, energii, dźwięczności i obwiedni widmowej. W zakresie poniżej 1 kbps nie można oczekiwać pełnej wierności względem sygnału oryginalnego. Bardziej realistycznym celem jest utrzymanie zrozumiałości oraz ograniczenie artefaktów, które szczególnie przeszkadzają w odbiorze mowy.

Obniżanie bitrate'u zwykle prowadzi do kilku rodzajów degradacji. Pierwszym z nich jest utrata szczegółów czasowych, widoczna między innymi w słabszym odwzorowaniu spółgłosek zwartych i szybkich przejść między fonemami. Drugim jest pogorszenie zgodności widmowej, które może objawiać się zmianą barwy, nadmiernym wygładzeniem lub zniekształceniem pasm istotnych dla rozróżniania głosek. Trzecim jest niestabilność rekonstrukcji, szczególnie w modelach neuronowych, gdzie dekoder może generować sygnał naturalnie brzmiący, ale częściowo niezgodny z treścią lub cechami mówcy. Z tego względu ocena jakości powinna obejmować nie tylko globalną miarę błędu, lecz także analizę zrozumiałości, artefaktów oraz stabilności słowników kodowych.

W projektowaniu kodeka mowy istotny jest więc kompromis między przepływnością, jakością i kosztem działania. Większy limit przepływności pozwala użyć większej liczby kodów lub częstszych ramek, co zwykle poprawia rekonstrukcję, ale oddala system od celu ultra-niskiej transmisji. Zbyt agresywne ograniczenie bitrate'u może natomiast doprowadzić do sytuacji, w której dekoder uzupełnia brakujące informacje na podstawie ogólnych wzorców z danych treningowych, a nie na podstawie rzeczywistego sygnału wejściowego. Dla urządzeń brzegowych dodatkowym ograniczeniem jest to, że poprawa jakości nie może wynikać wyłącznie z powiększenia modelu. Kodek musi zachować rozsądną przepływność, opóźnienie i koszt inferencji, dlatego dobór liczby ramek latentnych, słowników RVQ oraz rozmiaru słowników powinien być traktowany jako centralna decyzja projektowa.

\section{Podstawy neuronowej kompresji sygnału}
\label{sec:podstawy-neuronowej-kompresji}

Neuronowa kompresja sygnału audio polega na zastąpieniu części ręcznie projektowanych etapów kodeka modelem uczonym na danych. W klasycznych systemach kodowania znaczną część struktury narzuca projektant: wybiera transformację, parametry modelu mowy, sposób kwantyzacji i reguły rekonstrukcji. W podejściu neuronowym część tych decyzji zostaje przeniesiona do procesu uczenia, w którym model optymalizuje reprezentację pośrednią tak, aby umożliwiała rekonstrukcję sygnału przy ograniczonej liczbie przesyłanych informacji \cite{Goodfellow2016DeepLearning,Zeghidour2021SoundStream,Defossez2022Encodec}. Nie oznacza to jednak rezygnacji z wiedzy dziedzinowej. W praktyce architektura modelu nadal musi narzucać odpowiednie ograniczenia: małą częstotliwość ramek latentnych, skończony słownik kodów, ograniczoną liczbę warstw oraz możliwość uruchomienia na urządzeniu o małych zasobach.

Z punktu widzenia niniejszej pracy najważniejsze są te elementy uczenia głębokiego, które bezpośrednio wpływają na konstrukcję kodeka: autoenkoderowy podział na koder i dekoder, konwolucyjne przetwarzanie przebiegu czasowego, zwarta reprezentacja latentna, kwantyzacja wektorowa oraz dobór funkcji strat. Omówienie tych zagadnień ma charakter użytkowy. Celem nie jest pełny opis teorii sieci neuronowych, lecz wskazanie, dlaczego dana klasa mechanizmów jest przydatna przy kompresji mowy poniżej 1 kbps.

\subsection{Sieci neuronowe jako modele przetwarzania sygnałów}
\label{subsec:sieci-neuronowe-przetwarzanie-sygnalow}

Podstawową jednostką sieci neuronowej jest neuron obliczeniowy, który tworzy kombinację liniową wejść, dodaje przesunięcie, a następnie stosuje funkcję nieliniową. Historycznie jednym z pierwszych formalnych modeli tego typu był perceptron Rosenblatta \cite{Rosenblatt1958Perceptron}. Współczesne sieci głębokie składają się jednak z wielu warstw takich przekształceń, dzięki czemu mogą modelować zależności nieliniowe znacznie bardziej złożone niż pojedynczy klasyfikator liniowy \cite{Goodfellow2016DeepLearning}. Dla sygnałów audio jest to istotne, ponieważ zależność między przebiegiem czasowym, strukturą harmoniczną, artykulacją i percepcyjną jakością dźwięku nie jest prosta ani lokalnie liniowa.

Uczenie sieci polega na minimalizacji funkcji kosztu względem parametrów modelu. W praktyce wykorzystuje się algorytm propagacji wstecznej błędu, który umożliwia obliczanie gradientów dla kolejnych warstw modelu \cite{Rumelhart1986Backprop}. Parametry aktualizowane są następnie przez metodę optymalizacji, często w wariancie adaptacyjnym, takim jak Adam \cite{Kingma2015Adam}. W projektach implementacyjnych istotne są również biblioteki automatycznego różniczkowania, ponieważ pozwalają definiować złożone modele i funkcje strat bez ręcznego wyprowadzania pochodnych; przykładem takiego środowiska jest PyTorch \cite{Paszke2019PyTorch}. Z punktu widzenia kodeka mowy sam mechanizm uczenia jest jednak tylko narzędziem. Kluczowe pozostaje to, jakie ograniczenia zostaną narzucone reprezentacji pośredniej i czy model nauczy się przenosić przez wąski kanał informacje ważne percepcyjnie.

\subsection{Autoenkodery i wąskie gardło informacyjne}
\label{subsec:autoenkodery-waskie-gardlo}

Autoenkoder jest modelem złożonym z dwóch głównych części: kodera oraz dekodera. Koder przekształca sygnał wejściowy \(x\) do reprezentacji pośredniej \(z\), natomiast dekoder rekonstruuje sygnał \(\hat{x}\) na podstawie tej reprezentacji:

\begin{equation}
    z = f_{\theta}(x), \qquad \hat{x} = g_{\phi}(z).
\end{equation}

Parametry \(\theta\) i \(\phi\) są dobierane tak, aby zminimalizować błąd rekonstrukcji. Autoenkodery były szeroko wykorzystywane jako narzędzie uczenia reprezentacji i redukcji wymiarowości danych \cite{Hinton2006DimensionalityReduction,Goodfellow2016DeepLearning}. W kompresji ich rola jest szczególnie naturalna: koder odpowiada za wyodrębnienie informacji istotnej z punktu widzenia rekonstrukcji, a dekoder za odtworzenie sygnału z ograniczonego opisu.

Samo zastosowanie autoenkodera nie wystarcza jeszcze do zbudowania kodeka. Jeżeli reprezentacja \(z\) jest ciągła i ma dużą pojemność, może przechowywać bardzo dokładny opis wejścia, ale nie tworzy bezpośrednio skończonego strumienia bitów. Dlatego w systemach kompresji wprowadza się wąskie gardło informacyjne: zmniejsza się częstotliwość ramek latentnych, ogranicza wymiar wektorów, stosuje kwantyzację albo łączy te mechanizmy. W neuronowych kodekach audio, takich jak SoundStream i EnCodec, autoenkoder stanowi podstawę architektury, ale dopiero kwantyzacja wektorowa zamienia reprezentację latentną na dyskretny strumień indeksów \cite{Zeghidour2021SoundStream,Defossez2022Encodec}. W przypadku mowy poniżej 1 kbps wąskie gardło musi być bardzo silne, dlatego projekt kodera i kwantyzatora staje się ważniejszy niż sama zdolność dekodera do generowania naturalnego brzmienia.

\subsection{Konwolucje jednowymiarowe w analizie sygnału audio}
\label{subsec:konwolucje-jednowymiarowe-audio}

Konwolucje są jednym z podstawowych mechanizmów przetwarzania sygnałów przez sieci neuronowe. W ujęciu ogólnym warstwa konwolucyjna przesuwa niewielkie jądro po sygnale i stosuje ten sam zestaw wag w wielu położeniach. W sieciach przetwarzających obrazy konwolucje zostały spopularyzowane m.in. w pracach LeCuna i współautorów \cite{LeCun1998GradientBasedLearning}, natomiast w przypadku audio analogiczny mechanizm stosuje się wzdłuż osi czasu. Konwolucja jednowymiarowa może wykrywać lokalne wzorce w przebiegu, takie jak narastanie energii, struktury okresowe, fragmenty bezdźwięczne lub szybkie zmiany fazy. Współdzielenie wag jest korzystne, ponieważ te same lokalne zjawiska mogą występować w różnych momentach wypowiedzi.

W kodekach neuronowych konwolucje pełnią zwykle dwie funkcje. Po pierwsze, pozwalają budować lokalny opis sygnału, który jest odporniejszy na drobne przesunięcia w czasie niż analiza pojedynczych próbek. Po drugie, warstwy konwolucyjne z krokiem większym niż jeden umożliwiają redukcję długości sekwencji, a więc zmniejszenie liczby ramek, które trzeba później zakodować. W dekoderze analogiczną rolę pełni nadpróbkowanie, realizowane na przykład przez transponowane konwolucje lub przez powtarzanie próbek połączone z wygładzającą konwolucją. Przy projektowaniu kodeka strumieniowego istotne są również konwolucje przyczynowe, ponieważ ograniczają wykorzystanie informacji z przyszłości i zmniejszają opóźnienie.

W głębszych modelach konwolucyjnych często stosuje się połączenia rezydualne, które ułatwiają przepływ gradientu i pozwalają budować głębsze bloki bez tak silnej degradacji uczenia \cite{He2016DeepResidualLearning}. W audio podobny mechanizm może być używany w torze latentnym albo w dekoderze jako korekta rekonstrukcji. Należy jednak zachować ostrożność: większa liczba warstw poprawia pojemność modelu, ale jednocześnie zwiększa koszt inferencji, pamięć potrzebną na parametry oraz ryzyko przekroczenia ograniczeń urządzenia brzegowego.

\subsection{Reprezentacja latentna i redukcja czasowa}
\label{subsec:reprezentacja-latentna-redukcja-czasowa}

Reprezentacja latentna jest pośrednim opisem sygnału tworzonym przez koder. W przypadku kompresji audio ma ona zwykle postać sekwencji wektorów, a nie pojedynczego wektora dla całego nagrania. Jeżeli sygnał wejściowy ma długość \(T\) próbek, a koder zmniejsza rozdzielczość czasową \(S\)-krotnie, to liczba ramek latentnych wynosi w przybliżeniu \(T/S\). Dla sygnału próbkowanego z częstotliwością \(F_s\) częstotliwość ramek latentnych można zapisać jako:

\begin{equation}
    f_{\mathrm{lat}} = \frac{F_s}{S}.
\end{equation}

Ta zależność ma bezpośrednie znaczenie dla przepływności bitowej. Jeżeli koder wytwarza mniej ramek na sekundę, to kwantyzator musi przesłać mniej indeksów. Z drugiej strony zbyt silna redukcja czasowa może usunąć informacje potrzebne do wiernej rekonstrukcji spółgłosek, przejść między fonemami i szybkich zmian energii. W praktyce projektowanie reprezentacji latentnej wymaga więc kompromisu między pojemnością informacyjną, opóźnieniem i kosztem dekodowania.

Wymiar wektora latentnego wpływa na to, ile informacji może zostać przechowane w pojedynczej ramce przed kwantyzacją. Nie jest on jednak równoważny przepływności bitowej, ponieważ w kodeku dyskretnym transmitowany jest indeks kodu, a nie pełny wektor zmiennoprzecinkowy. Można więc zwiększać wymiar latentny, aby ułatwić pracę koderowi i dekoderowi, zachowując tę samą nominalną liczbę bitów, o ile liczba słowników i liczba kodów w słownikach pozostają bez zmian. Takie rozróżnienie jest ważne w neuronowych kodekach end-to-end, ponieważ jakość rekonstrukcji zależy zarówno od geometrii przestrzeni latentnej, jak i od dyskretnej pojemności kwantyzatora \cite{Zeghidour2021SoundStream,Defossez2022Encodec}.

\subsection{Kwantyzacja wektorowa i rezydualna kwantyzacja wektorowa}
\label{subsec:kwantyzacja-wektorowa-rvq}

Kwantyzacja wektorowa polega na zastąpieniu ciągłego wektora jednym z elementów skończonego słownika kodowego. Dla słownika \(\mathcal{C} = \{e_1, e_2, \ldots, e_K\}\) wektor latentny jest przypisywany do najbliższego kodu, a do strumienia bitów trafia tylko indeks wybranego elementu. Klasyczna teoria kwantyzacji wektorowej była rozwijana znacznie wcześniej niż współczesne kodeki neuronowe \cite{Gray1984VectorQuantization}, ale w modelach głębokich problem ten pojawia się ponownie jako sposób uczenia dyskretnych reprezentacji. VQ-VAE pokazał, że kwantyzację można włączyć do autoenkodera i trenować reprezentacje dyskretne z użyciem estymatora prostego przejścia gradientu oraz dodatkowych składników stabilizujących słownik \cite{VanDenOord2017NeuralDiscrete}.

Jeżeli pojedynczy słownik ma \(K\) kodów, to jeden indeks wymaga nominalnie \(\log_2 K\) bitów. Dla wielu słowników, stosowanych w każdej ramce latentnej, przepływność indeksów można zapisać jako:

\begin{equation}
    R = f_{\mathrm{lat}} \sum_{i=1}^{N_q} \log_2 K_i,
\end{equation}

gdzie \(N_q\) oznacza liczbę etapów kwantyzacji, a \(K_i\) liczbę kodów w \(i\)-tym słowniku. W tym zapisie \(\log_2 K_i\) oznacza nominalną informację potrzebną do wyboru kodu ze słownika. Przy stałoszerokim pakowaniu indeksów należy użyć \(\lceil \log_2 K_i \rceil\), jeśli rozmiar słownika nie jest potęgą dwójki. W analizowanym wariancie \(K_i=32\), dlatego oba zapisy prowadzą do \(5\) bitów na indeks. Widać z tego, że przepływność zależy bezpośrednio od częstotliwości ramek latentnych, liczby słowników oraz ich rozmiaru. Przy projektowaniu kodeka poniżej 1 kbps nawet niewielka zmiana tych parametrów może mieć duże znaczenie, ponieważ liczba dostępnych bitów jest bardzo mała.

Rezydualna kwantyzacja wektorowa rozwija ten mechanizm przez użycie wielu etapów. Pierwszy słownik koduje przybliżenie wektora, kolejny koduje błąd pozostały po pierwszym etapie, następny koduje kolejną resztę i tak dalej. Dzięki temu można zwiększyć dokładność rekonstrukcji latentów bez gwałtownego zwiększania rozmiaru pojedynczego słownika. Taki mechanizm jest wykorzystywany w kodekach SoundStream i EnCodec, a warianty grupowej rezydualnej kwantyzacji pojawiają się również w HiFi-Codec \cite{Zeghidour2021SoundStream,Defossez2022Encodec,Yang2023HiFiCodec}. Ograniczeniem RVQ jest możliwość nierównego wykorzystania kodów, czyli sytuacja, w której model wybiera tylko niewielką część słownika. Z tego powodu podczas treningu stosuje się dodatkowe straty, regularyzację wykorzystania kodów albo procedury przeciwdziałające zapadaniu się słowników.

\begin{center}
\centering
\footnotesize
\setlength{\tabcolsep}{2pt}
\begin{tabular}{c c c c c c c}
\fbox{\parbox{0.11\textwidth}{\centering latent\\\(z\)}} &
\(\rightarrow\) &
\fbox{\parbox{0.12\textwidth}{\centering słownik\\\(C_1\)}} &
\(\rightarrow\) &
\fbox{\parbox{0.12\textwidth}{\centering reszta\\\(r_1\)}} &
\(\rightarrow\) &
\fbox{\parbox{0.12\textwidth}{\centering słownik\\\(C_2\)}} \\
 & & indeks \(i_1\) & & & & indeks \(i_2\) \\
\end{tabular}

\vspace{0.4em}

\begin{tabular}{c c c c c}
\fbox{\parbox{0.12\textwidth}{\centering kolejne\\reszty}} &
\(\rightarrow\) &
\fbox{\parbox{0.16\textwidth}{\centering indeksy\\\(i_1,\ldots,i_N\)}} &
\(\rightarrow\) &
\fbox{\parbox{0.16\textwidth}{\centering suma\\kodów}} \\
\end{tabular}
\captionof{figure}{Idea rezydualnej kwantyzacji wektorowej: kolejne słowniki kodują pozostały błąd reprezentacji.}
\label{fig:rvq-reszty}
\end{center}

\subsection{Funkcje aktywacji, normalizacja i regularyzacja}
\label{subsec:aktywacje-normalizacja-regularyzacja}

Funkcje aktywacji wprowadzają nieliniowość między kolejnymi warstwami sieci. Bez nich złożenie wielu warstw liniowych nadal byłoby przekształceniem liniowym, a model nie mógłby opisywać złożonych zależności obecnych w sygnale mowy. W praktyce stosuje się różne funkcje aktywacji, takie jak ReLU, GELU, sigmoid, tangens hiperboliczny lub aktywacje dostosowane do struktur okresowych. Wybór aktywacji wpływa na stabilność gradientów, zakres wartości pośrednich oraz zdolność modelu do odtwarzania drobnych szczegółów sygnału \cite{Goodfellow2016DeepLearning}. W kodeku audio nie jest to wyłącznie decyzja numeryczna: zbyt agresywna nieliniowość może prowadzić do zniekształceń, a zbyt łagodna może ograniczyć zdolność modelu do rekonstrukcji szybkich zmian w przebiegu.

Normalizacja służy stabilizacji uczenia przez kontrolowanie rozkładu aktywacji. W przetwarzaniu sekwencji często stosuje się normalizację warstwową, która oblicza statystyki dla pojedynczego przykładu, a nie dla całego mini-batcha \cite{Ba2016LayerNormalization}. Może to być korzystne w modelach audio, ponieważ długość sygnałów, dynamika głośności i zróżnicowanie mówców powodują duże zmiany skali cech. Jednocześnie normalizacja umieszczona bezpośrednio przed kwantyzacją może zmieniać geometrię przestrzeni latentnej, dlatego jej wpływ na stabilność RVQ wymaga weryfikacji eksperymentalnej.

Regularyzacja obejmuje techniki ograniczające przeuczenie i niestabilność modelu. Klasycznym przykładem jest dropout, który losowo wyłącza część aktywacji podczas treningu \cite{Srivastava2014Dropout}. W kodekach audio stosowanie takiej regularyzacji wymaga jednak ostrożności, ponieważ celem nie jest klasyfikacja, lecz precyzyjna rekonstrukcja sygnału. Naturalną formą regularyzacji jest również samo wąskie gardło kompresyjne: ograniczona liczba kodów, ograniczona liczba ramek latentnych oraz kary związane z kwantyzacją zmuszają model do wyboru informacji najważniejszych dla rekonstrukcji. W tym sensie architektura kodeka pełni funkcję regularyzującą równie istotną jak klasyczne techniki uczenia głębokiego.

\subsection{Funkcje strat stosowane w rekonstrukcji audio}
\label{subsec:funkcje-strat-rekonstrukcja-audio}

Funkcja straty określa, jakie błędy model ma uznawać za najważniejsze podczas treningu. W najprostszym przypadku można minimalizować różnicę między przebiegiem oryginalnym i zrekonstruowanym bezpośrednio w dziedzinie czasu, na przykład za pomocą normy \(L_1\) lub \(L_2\). Takie podejście dobrze wiąże fazę i amplitudę próbek, ale bywa niewystarczające percepcyjnie. Dwa sygnały o podobnym brzmieniu mogą różnić się punktowo w czasie, a z kolei niski błąd próbkowy nie gwarantuje zachowania struktury harmonicznej i barwy.

Z tego powodu w modelach generujących lub rekonstruujących mowę często stosuje się straty widmowe. Porównują one sygnały po transformacji STFT, dzięki czemu model otrzymuje informację o błędach w dziedzinie czasu i częstotliwości. Szczególnie użyteczne są straty wielorozdzielcze, w których porównanie wykonywane jest dla kilku długości okna i kilku kroków analizy. Takie podejście zostało wykorzystane m.in. w Parallel WaveGAN, gdzie wielorozdzielcza strata spektrogramowa wspierała generowanie fali dźwiękowej \cite{Yamamoto2020ParallelWaveGAN}. W kodekach neuronowych podobną rolę pełnią straty rekonstrukcyjne, widmowe, percepcyjne i adwersarialne, łączone w celu poprawy naturalności odtwarzanego sygnału \cite{Zeghidour2021SoundStream,Defossez2022Encodec}.

W systemie z kwantyzacją konieczne są również składniki związane z samym słownikiem kodowym. Strata przywiązania do kodu, określana w literaturze jako \emph{commitment loss}, ogranicza odrywanie się reprezentacji latentnej od wybranych kodów, natomiast składniki promujące wykorzystanie słownika przeciwdziałają sytuacji, w której większość kodów pozostaje nieużywana. Dodatkowe straty, takie jak podobieństwo kosinusowe przebiegów, straty melowe lub kary za nadmiarową energię w wybranych pasmach, mogą poprawiać wybrane aspekty rekonstrukcji, ale każda z nich wprowadza własny kompromis. Zbyt duża waga strat widmowych może prowadzić do sztucznego wygładzenia lub wzrostu szumu tła, natomiast zbyt silne wymuszanie zgodności w czasie może pogarszać naturalność brzmienia. Dlatego dobór funkcji strat w neuronowym kodeku mowy jest częścią projektu systemu, a nie jedynie technicznym detalem treningu.

\section{Metody kodowania dźwięku przy bardzo niskich przepływnościach}
\label{sec:metody-kodowania-ultra-low}

Metody kodowania przy bardzo niskich przepływnościach różnią się przede wszystkim tym, jak silne założenia przyjmują o strukturze sygnału. Przy klasycznej kompresji audio można zwykle przesyłać stosunkowo bogaty opis widmowy, a zadaniem kodeka jest usunięcie informacji nadmiarowej lub mniej istotnej percepcyjnie. W zakresie poniżej 1 kbps sytuacja jest inna: liczba bitów jest tak mała, że kodek musi zdecydować, które informacje są konieczne dla zachowania zrozumiałości i podstawowej naturalności mowy. Dlatego w praktyce stosuje się podejścia, które nie próbują zachować pełnego przebiegu akustycznego, lecz opisują go przez parametry, predykcję, reprezentację latentną albo symbole semantyczne.

\subsection{Podejścia parametryczne, predykcyjne i transformacyjne}
\label{subsec:parametryczne-predykcyjne-transformacyjne}

Najbardziej bezpośrednią drogą do bardzo niskiego bitrate'u jest kodowanie parametryczne. Kodek nie przesyła wtedy próbek ani pełnego widma, lecz niewielki zestaw parametrów opisujących model mowy, zwykle obejmujący częstotliwość podstawową, energię, dźwięczność i obwiednię widmową. Takie podejście jest zgodne z modelem źródło-filtr: zamiast kodować całą falę, opisuje się pobudzenie i filtr traktowany jako uproszczony model toru głosowego. Codec 2 oraz MELPe pokazują, że dzięki temu można uzyskać przepływności subkilobitowe, w tym zakresy rzędu 700, 600 lub nawet 450 bit/s \cite{Rowe2019Codec2,Rowe2021Codec2450,NATO2015STANAG4591}. Ograniczeniem tego kierunku pozostaje jednak naturalność rekonstrukcji, ponieważ model parametryczny upraszcza wiele szczegółów artykulacyjnych i akustycznych.

Metody predykcyjne wykorzystują zależność między kolejnymi próbkami lub parametrami sygnału. W kodowaniu mowy szczególnie ważna jest liniowa predykcja, która pozwala modelować obwiednię widmową krótkiego fragmentu i oddzielać ją od pobudzenia. Bäckström i Mägi wskazują, że predykcja rzadka może poprawiać efektywność kodowania przy małych przepływnościach, ponieważ ogranicza liczbę istotnych współczynników modelu \cite{Backstrom2017SparseLinearPrediction}. W praktyce metody predykcyjne dobrze pasują do mowy, lecz gorzej radzą sobie z fragmentami nietypowymi, szumowymi lub silnie zmiennymi.

Kodowanie transformacyjne opiera się na zamianie sygnału na reprezentację częstotliwościową i kwantyzacji współczynników transformaty. Przy wyższych bitrate'ach pozwala ono wykorzystywać właściwości słuchu, na przykład przez dokładniejsze kodowanie pasm istotnych percepcyjnie. W zakresie subkilobitowym sam opis transformacyjny staje się jednak zbyt kosztowny, szczególnie dla sygnału mowy 16 kHz, ponieważ nawet silnie zredukowany spektrogram zawiera wiele wartości na sekundę. Reprezentacje STFT i melowe pozostają natomiast użyteczne jako narzędzia analizy, funkcje strat i metryki diagnostyczne \cite{RabinerSchafer2011SpeechDSP}.

\subsection{Podejścia hybrydowe i neuronowe}
\label{subsec:podejscia-hybrydowe-neuronowe}

Podejścia hybrydowe łączą klasyczny model mowy z komponentem neuronowym. Ich celem nie jest całkowite zastąpienie wiedzy dziedzinowej siecią neuronową, lecz zmniejszenie trudności zadania rekonstrukcji. LPCNet wykorzystuje predykcję liniową, a sieć neuronowa wspiera generowanie sygnału mowy na podstawie ograniczonego opisu \cite{Valin2019LPCNet}. Prace nad optymalizacją LPCNet dla urządzeń embedded pokazują, że taki kierunek jest istotny dla przetwarzania brzegowego, ponieważ jakość modelu musi być rozpatrywana razem z kosztem obliczeniowym i możliwością pracy w czasie rzeczywistym \cite{Zhu2021OptimizingLPCNet,Valin2022UltraLowLatencySpeechEnhancement}. Ograniczeniem rozwiązań hybrydowych jest zależność od jawnie wybranego modelu mowy oraz konieczność bardzo silnej redukcji parametrów przy zejściu poniżej 1 kbps.

Neuronowe kodeki end-to-end przesuwają większą część projektu do procesu uczenia. Koder tworzy reprezentację latentną, kwantyzer zamienia ją na dyskretne indeksy, a dekoder rekonstruuje przebieg audio. SoundStream oraz EnCodec pokazują, że taki układ może osiągać wysoką jakość przy małej lub umiarkowanej przepływności, zwłaszcza gdy wykorzystuje się rezydualną kwantyzację wektorową \cite{Zeghidour2021SoundStream,Defossez2022Encodec}. HiFi-Codec rozwija ten kierunek przez grupową rezydualną kwantyzację \cite{Yang2023HiFiCodec}. W kontekście niniejszej pracy najważniejsza jest nie sama jakość tych systemów, lecz mechanizm dyskretnego wąskiego gardła: fala audio może być opisywana sekwencją indeksów, a nie bezpośrednio współczynnikami klasycznego modelu mowy.

Przeniesienie tych metod do zakresu poniżej 1 kbps wymaga zmniejszenia liczby ramek latentnych, rozmiaru słowników lub liczby etapów kwantyzacji, co zwiększa błąd kwantyzacji i obciąża dekoder. Prace oparte na WaveNet lub na uczeniu end-to-end wskazują, że generatywne modele fali mogą poprawiać naturalność rekonstrukcji, lecz zwykle są kosztowne obliczeniowo \cite{Kleijn2018WaveNetLowRate,Kankanahalli2018EndToEndSpeechCoding,VanDenOord2016WaveNet}. Z kolei badania nad efektywnymi neuronowymi kodekami dla strumieniowania pokazują, że opóźnienie i koszt inferencji są osobnymi ograniczeniami projektowymi \cite{Morrison2022EfficientNeuralAudioCodecs}.

\subsection{Reprezentacje semantyczne i wybór kierunku w projekcie}
\label{subsec:reprezentacje-semantyczne-wybor-kierunku}

Osobną grupę stanowią metody oparte na reprezentacjach semantycznych lub samonadzorowanych. Zamiast przesyłać pełny opis akustyczny, kodują one informację bliższą treści wypowiedzi, fonemom lub jednostkom dyskretnym wyznaczonym przez model uczony na dużych korpusach mowy. Speech resynthesis z dyskretnych reprezentacji samonadzorowanych oraz SpeechTokenizer pokazują, że możliwa jest rekonstrukcja mowy z tokenów częściowo oddzielających treść od cech akustycznych \cite{Polyak2021SpeechResynthesis,Zhang2023SpeechTokenizer}. Modele takie jak HuBERT i ContentVec mogą być użyteczne przy ocenie zgodności treści \cite{Hsu2021HuBERT,Qian2022ContentVec}. Dla kodeka subkilobitowego jest to atrakcyjny kierunek, ponieważ treść językowa może być znacznie bardziej zwarta niż pełna fala audio.

Podejścia semantyczne mają jednak istotne ograniczenia. Ekstrakcja cech z dużych modeli samonadzorowanych może być kosztowna, a reprezentacje semantyczne często zachowują treść lepiej niż barwę, prozodię i szczegóły artykulacyjne. Rekonstrukcja z tokenów może więc prowadzić do sygnału poprawnego językowo, ale mniej wiernego względem oryginalnego mówcy lub emocji wypowiedzi. Z tego względu w praktycznym kodeku mowy poniżej 1 kbps podejście semantyczne może być używane pomocniczo, na przykład w funkcji straty lub analizie zgodności.

W projekcie \texttt{sirencodec} przyjęto rozwiązanie pośrednie między klasycznym kodowaniem parametrycznym a pełnym kodekiem neuronowym wysokiej przepływności. Z metod parametrycznych zaczerpnięto założenie, że przy ultra-niskim bitrate'cie nie należy dążyć do pełnej zgodności próbek, lecz do zachowania informacji najważniejszej dla mowy. Z kodeków neuronowych wykorzystano autoenkoderowy podział na koder i dekoder, redukcję czasową oraz rezydualną kwantyzację wektorową. Z podejść transformacyjnych wykorzystano reprezentacje STFT i melowe jako narzędzia uczenia i oceny, a nie jako bezpośredni strumień transmisji. Nie przyjęto natomiast ciężkiego modelu semantycznego jako obowiązkowego elementu inferencji, ponieważ zwiększałoby to koszt wdrożenia brzegowego. Taki wybór odpowiada głównemu ograniczeniu pracy: system ma działać poniżej 1 kbps, ale jednocześnie powinien pozostać możliwy do analizy, implementacji i uruchomienia w środowisku o ograniczonych zasobach.

\begin{center}
\footnotesize
\setlength{\tabcolsep}{4pt}
\begin{tabularx}{0.96\textwidth}{>{\raggedright\arraybackslash}p{0.20\textwidth} >{\raggedright\arraybackslash}X >{\raggedright\arraybackslash}X >{\raggedright\arraybackslash}X}
\toprule
\textbf{Rodzina metod} & \textbf{Główna idea} & \textbf{Zaleta przy niskim bitrate} & \textbf{Ograniczenie} \\
\midrule
Parametryczne & Opis mowy przez parametry źródła i filtra. & Bardzo mała liczba bitów i niski koszt obliczeniowy. & Ograniczona naturalność i silne uproszczenie sygnału. \\
Predykcyjne & Wykorzystanie zależności czasowych i modelu LPC. & Dobra zgodność z lokalną strukturą mowy. & Wrażliwość na błędy modelu i nietypowe fragmenty. \\
Transformacyjne & Kwantyzacja reprezentacji widmowej. & Dobra kontrola pasm istotnych percepcyjnie. & Zbyt duży opis dla zakresu subkilobitowego. \\
Neuronowe autoenkodery & Uczenie zwartej reprezentacji latentnej i dekodera fali. & Elastyczność i możliwość optymalizacji end-to-end. & Większy koszt treningu, inferencji i ryzyko niestabilności. \\
Kwantyzacja wektorowa & Zamiana latentów na indeksy ze słowników. & Jasna kontrola nominalnej przepływności. & Możliwe zapadanie słowników i błąd kwantyzacji. \\
Semantyczne & Kodowanie treści lub tokenów mowy. & Potencjalnie bardzo zwarty opis informacji językowej. & Ryzyko utraty barwy, prozodii i zgodności z mówcą. \\
\bottomrule
\end{tabularx}
\captionof{table}{Porównanie rodzin metod stosowanych lub rozważanych przy bardzo niskich przepływnościach.}
\label{tab:metody-ultra-low-bitrate}
\end{center}

\section{Metryki oceny jakości rekonstrukcji dźwięku}
\label{sec:metryki-jakosci}

Ocena jakości rekonstrukcji dźwięku wymaga użycia kilku uzupełniających się miar. W kodeku mowy nie wystarcza pojedyncza liczba, ponieważ różne metryki opisują różne aspekty sygnału: zgodność próbek, strukturę widmową, zrozumiałość wypowiedzi albo przewidywaną jakość percepcyjną. Jest to szczególnie ważne przy bardzo niskich przepływnościach, gdzie rekonstrukcja może być zrozumiała, ale mniej wierna próbkowo, albo odwrotnie: może zachowywać część energii sygnału, lecz zawierać artefakty utrudniające odbiór. Z tego powodu w pracy przyjęto zestaw metryk, a nie jedną miarę rozstrzygającą.

\subsection{Miary zgodności w dziedzinie czasu}
\label{subsec:metryki-czasowe}

Najprostszą grupę stanowią miary obliczane bezpośrednio na przebiegu czasowym. Błąd \(L_1\) mierzy średnią wartość bezwzględną różnicy między sygnałem referencyjnym \(x\) i rekonstrukcją \(\hat{x}\):

\begin{equation}
    L_1(x,\hat{x}) = \frac{1}{T} \sum_{t=1}^{T} |x_t - \hat{x}_t|.
\end{equation}

Miara ta jest łatwa do interpretacji i przydatna podczas treningu, ponieważ bezpośrednio karze różnice amplitudy. Jej ograniczeniem jest słaba zgodność z percepcją. Niewielkie przesunięcie fazowe, lokalne przesunięcie w czasie albo zmiana amplitudy mogą znacząco zwiększyć błąd \(L_1\), nawet jeżeli sygnał nadal brzmi podobnie. Z drugiej strony niski błąd próbkowy nie gwarantuje naturalności, ponieważ może nie ujawniać zniekształceń widmowych istotnych dla barwy mowy.

Uzupełnieniem jest podobieństwo kosinusowe, które mierzy zgodność kierunku dwóch wektorów sygnału, a nie bezwzględną różnicę amplitud:

\begin{equation}
    \cos(x,\hat{x}) = \frac{\langle x,\hat{x}\rangle}{\|x\|_2 \|\hat{x}\|_2}.
\end{equation}

Wartość bliska \(1\) oznacza dużą zgodność kształtu przebiegu po uwzględnieniu skali. Metryka ta bywa użyteczna przy porównywaniu rekonstrukcji, ponieważ jest mniej wrażliwa na globalne przeskalowanie amplitudy niż \(L_1\). Nie rozwiązuje jednak problemu percepcji: dwa sygnały o wysokim podobieństwie kosinusowym mogą nadal różnić się lokalnymi artefaktami, a sygnały podobne odsłuchowo mogą otrzymać gorszy wynik przy przesunięciu czasowym.

Bardziej wyspecjalizowaną miarą jest SI-SDR, czyli \emph{scale-invariant signal-to-distortion ratio}. Metryka ta porównuje estymowany sygnał z referencją po dopuszczeniu optymalnego przeskalowania, dzięki czemu nie karze samej różnicy poziomu głośności. Dla referencji \(x\) i rekonstrukcji \(\hat{x}\) wyznacza się składową zgodną z referencją:

\begin{equation}
    x_{\mathrm{target}} = \frac{\langle \hat{x}, x\rangle}{\|x\|_2^2}x,
    \qquad
    e_{\mathrm{noise}} = \hat{x} - x_{\mathrm{target}},
\end{equation}

a następnie:

\begin{equation}
    \mathrm{SI\mbox{-}SDR} = 10\log_{10}\frac{\|x_{\mathrm{target}}\|_2^2}{\|e_{\mathrm{noise}}\|_2^2}.
\end{equation}

Le Roux i współautorzy wskazali, że klasyczne warianty SDR mogą być mylące w niektórych zadaniach separacji i rekonstrukcji sygnału, a wersja skalo-niezmiennicza lepiej kontroluje wpływ globalnej amplitudy \cite{LeRoux2019SDR}. W implementacji wykorzystywanej w pracy sygnały są przed obliczeniem SI-SDR przycinane do wspólnej długości i centrowane przez odjęcie wartości średniej. W kodeku mowy SI-SDR jest użytecznym wskaźnikiem zgodności sygnału w dziedzinie czasu, lecz nadal nie powinien być interpretowany jako pełna miara jakości odsłuchowej. Przy silnej kompresji model może uzyskać ograniczony SI-SDR, mimo że zrekonstruowana mowa pozostaje zrozumiała.

\subsection{Miary widmowe}
\label{subsec:metryki-widmowe}

Miary widmowe porównują sygnały po transformacji do dziedziny częstotliwości. Są szczególnie ważne w ocenie kodeków mowy, ponieważ wiele artefaktów jest lepiej widocznych w spektrogramie niż w przebiegu czasowym. Logarytmiczna odległość widmowa, określana jako LSD, mierzy różnicę między widmami w skali logarytmicznej. Dla pojedynczej ramki można ją zapisać przykładowo jako:

\begin{equation}
    \mathrm{LSD} =
    \sqrt{\frac{1}{K}\sum_{k=1}^{K}
    \left(20\log_{10}(|X_k|+\epsilon) -
    20\log_{10}(|\hat{X}_k|+\epsilon)\right)^2},
\end{equation}

gdzie \(X_k\) i \(\hat{X}_k\) oznaczają widma sygnału referencyjnego oraz zrekonstruowanego, a \(\epsilon\) zabezpiecza obliczenia przed logarytmem zera. Miary odległości widmowej były stosowane w przetwarzaniu mowy już w klasycznych pracach dotyczących porównywania reprezentacji spektralnych \cite{GrayMarkel1976DistanceMeasures}. W kontekście kodeka neuronowego LSD pomaga wykrywać utratę energii w pasmach częstotliwości, zmianę obwiedni widmowej i nadmierne wygładzenie rekonstrukcji.

Warto jednak zaznaczyć, że LSD nie jest jednoznaczną miarą jakości percepcyjnej. Może silnie reagować na różnice w pasmach, które nie są jednakowo ważne dla odbioru mowy, a sposób obliczania zależy od parametrów STFT, liczby pasm, uśredniania ramek i stabilizacji logarytmu. Dlatego w analizie wyników miara ta powinna być zestawiana z metrykami percepcyjnymi oraz odsłuchem. Jej zaletą jest natomiast przejrzystość diagnostyczna: wysoka odległość widmowa zwykle wskazuje, że model nie odtworzył poprawnie struktury częstotliwościowej sygnału.

\subsection{Miary percepcyjne i zrozumiałości}
\label{subsec:metryki-percepcyjne-zrozumialosc}

PESQ jest jedną z klasycznych obiektywnych miar jakości mowy. Została opisana jako metoda porównująca sygnał referencyjny i zdegradowany z uwzględnieniem modelu percepcyjnego, a następnie znormalizowana w rekomendacji ITU-T P.862 dla oceny jakości w sieciach telefonicznych i kodekach mowy \cite{Rix2001PESQ,ITUT2001P862}. Nowszym następcą w rodzinie standardów ITU jest POLQA, ujęta w rekomendacji P.863 \cite{ITUT2018P863}. Wynik PESQ jest zwykle interpretowany jako przybliżenie jakości odsłuchowej, dlatego bywa wygodny przy porównywaniu wariantów kodeka. Ograniczeniem jest jednak zakres warunków, dla których metryka była projektowana. PESQ dobrze pasuje do wielu degradacji telekomunikacyjnych, ale może nie w pełni odzwierciedlać artefakty generowane przez neuronowe dekodery, zwłaszcza gdy rekonstrukcja jest naturalna, lecz częściowo niezgodna z sygnałem referencyjnym.

STOI, czyli \emph{short-time objective intelligibility}, jest miarą nastawioną na przewidywanie zrozumiałości mowy, a nie ogólnej naturalności brzmienia. Taal i współautorzy zaproponowali ją dla sygnałów mowy zniekształcanych m.in. przez szum i przetwarzanie czasowo-częstotliwościowe \cite{Taal2011STOI}. Z punktu widzenia niniejszej pracy jest to metryka istotna, ponieważ przy przepływności poniżej 1 kbps zachowanie zrozumiałości jest ważniejsze niż pełna wierność akustyczna. Wysoki wynik STOI może sugerować, że treść wypowiedzi została zachowana, ale nie oznacza automatycznie naturalnego głosu, stabilnej barwy ani braku artefaktów.

ViSQOL jest pełnoreferencyjną metryką jakości, która porównuje sygnał referencyjny i testowy przez podobieństwo spektro-czasowe, a wynik może być mapowany na skalę MOS-LQO \cite{Hines2015ViSQOL}. W odróżnieniu od prostych miar błędu, ViSQOL próbuje lepiej odzwierciedlać subiektywną jakość odsłuchową. Jest to przydatne w ocenie kodeków, ponieważ pozwala porównywać warianty systemu w sposób bliższy testom percepcyjnym, choć nadal automatyczny. Podobnie jak PESQ, metryka ta ma jednak ograniczenia wynikające z danych i typów zniekształceń, na których była rozwijana. Wynik MOS-LQO należy więc traktować jako wskaźnik pomocniczy, a nie równoważnik formalnego testu odsłuchowego.

\begin{center}
\footnotesize
\setlength{\tabcolsep}{4pt}
\begin{tabularx}{0.96\textwidth}{>{\raggedright\arraybackslash}p{0.16\textwidth} >{\raggedright\arraybackslash}p{0.22\textwidth} >{\raggedright\arraybackslash}X >{\raggedright\arraybackslash}X}
\toprule
\textbf{Metryka} & \textbf{Zakres oceny} & \textbf{Interpretacja} & \textbf{Główne ograniczenie} \\
\midrule
\(L_1\) & Dziedzina czasu & Średnia różnica amplitudy próbek. & Słaba zgodność z percepcją i wrażliwość na przesunięcia. \\
Cos & Dziedzina czasu & Zgodność kształtu przebiegu po normalizacji skali. & Nie wykrywa wszystkich artefaktów lokalnych i widmowych. \\
SI-SDR & Dziedzina czasu & Stosunek sygnału zgodnego z referencją do zniekształceń. & Nie opisuje bezpośrednio naturalności ani zrozumiałości. \\
LSD & Dziedzina częstotliwości & Różnica między logarytmicznymi widmami. & Zależy od parametrów analizy i nie waży pasm percepcyjnie. \\
PESQ & Jakość mowy & Predykcja jakości odsłuchowej względem referencji. & Ograniczona zgodność dla nietypowych artefaktów neuronowych. \\
STOI & Zrozumiałość mowy & Przewidywana zrozumiałość sygnału mowy. & Nie mierzy naturalności głosu ani barwy. \\
ViSQOL & Jakość percepcyjna & MOS-LQO na podstawie podobieństwa spektro-czasowego. & Nadal pozostaje estymatorem, a nie testem odsłuchowym. \\
\bottomrule
\end{tabularx}
\captionof{table}{Metryki wykorzystywane do oceny rekonstrukcji dźwięku w pracy.}
\label{tab:metryki-jakosci}
\end{center}

\subsection{Interpretacja wyników w kodeku ultra-niskobitowym}
\label{subsec:interpretacja-metryk-ultra-low}

Przy ocenie kodeka o ultra-niskiej przepływności najważniejsze jest zestawianie metryk, a nie wybór jednej wartości granicznej. Wysoki SI-SDR lub niski błąd \(L_1\) może wskazywać na lepszą zgodność przebiegu, ale nie musi oznaczać lepszego odbioru mowy. Z kolei wyższy STOI może potwierdzać zachowanie treści wypowiedzi, lecz nie mówi wystarczająco dużo o barwie, szumie tła i naturalności. PESQ oraz ViSQOL są bliższe jakości percepcyjnej, ale również mogą być mylące, jeżeli artefakty kodeka różnią się od zniekształceń typowych dla klasycznych systemów telekomunikacyjnych.

Z tego powodu w analizie wyników należy traktować metryki jako narzędzia diagnostyczne. Jeżeli kilka miar poprawia się jednocześnie, można ostrożnie wnioskować o lepszej rekonstrukcji. Jeżeli metryki są sprzeczne, konieczna jest interpretacja jakościowa, obejmująca spektrogramy oraz odsłuch przykładów. W niniejszej pracy szczególne znaczenie ma relacja między zrozumiałością a naturalnością: przy przepływności poniżej 1 kbps akceptowalna rekonstrukcja może mieć ograniczoną wierność akustyczną, ale powinna zachowywać treść wypowiedzi, rytm i podstawowe cechy głosu. Automatyczne metryki ułatwiają porównywanie wielu wariantów modelu, lecz końcowa ocena wymaga uwzględnienia ograniczeń każdej z nich.
